\chapter{Models of Limit Order Books}
\label{chap:lob_models}

This chapter is dedicated to a review of agent-based models of limit order books, that is,
models depicting, at the individual agent level, possibly from a statistical point of view, the
interactions that lead to transactions on a financial market. Far from being exhaustive,
this survey is based on selected models that are representative of important,
specific trends in agent-based modelling.

\section{Early Order-driven Market Modelling: Market Microstructure and Policy}
The earliest market simulations were designed to study regulatory issues. The first
attempt, by Stigler \cite{stigler1964}, investigated the impact of SEC regulations on U.S.\ stock
markets using historical data from the 1920s and 1950s. With the growing computerization
of markets, Hakansson et al.\ \cite{hakansson1985} developed a simulator to test the feasibility of
automated market making.

While we do not review the full microstructure literature (see O'Hara \cite{ohara1997} and
Hasbrouck \cite{hasbrouck2007}), we highlight a few foundational works that underpin modern
order book modelling.

\subsubsection{A Pioneer Order Book Model}
Stigler's 1964 model is the earliest documented attempt to simulate a financial market.
Motivated by criticisms of the SEC’s 1963 Special Study \cite{cohen1963b,cohen1963a}, Stigler
proposes a simple simulation in which price evolves within a discrete grid of $L=10$ ticks.
Orders arrive randomly, namely, buy or sell with equal probability and uniformly over the price
grid, and become market orders if they cross the best opposing quote. Orders not executed
within $N=25$ time steps are cancelled. All orders have unit size, and the maximum number
of standing orders is $N$.

Slanina \cite{slanina2008} later simulated a similar model with $L=5000$ and $N=5000$,
showing that returns follow a heavy-tailed distribution with a power-law exponent of about
0.3, though limited by the artificial cutoff imposed by $L$.

\subsubsection{Microstructure of the Double Auction}
Garman \cite{garman1976} provides an early study of the continuous double auction that explicitly
incorporates temporal order flow. Buy and sell orders arrive as independent Poisson
processes with intensities $\lambda$ and $\mu$, respectively, and are of unit size. Orders crossing
the best quote become market orders. The state of the book is represented as a vector of
standing quantities at each price level. Although analytically tractable in principle, the
model quickly becomes complex, highlighting the need for simulation methods and
numerical analysis.

The key question posed by Garman is structural: Does the double auction mechanism
itself generate the leptokurtosis observed in empirical return series.

\subsubsection{Zero-intelligence Models}
The Stigler and Garman models belong to a broad class of ``zero-intelligence'' models,
later formalized by Gode and Sunder \cite{gode1993}. In such models, orders are submitted
randomly, without strategic behavior. Two variants are typically considered:
\begin{itemize}
    \item \emph{Unconstrained zero-intelligence traders}: random orders at random prices.
    \item \emph{Constrained zero-intelligence traders}: random orders respecting a personal reference
    price, preventing trades at a loss.
\end{itemize}

Zero-intelligence markets exhibit surprising allocative efficiency but also unrealistic price
volatility. Later critiques (e.g.\ \cite{cliffbruten1997}) argue that some of these results are artefacts of
the imposed constraints.

These early models formed the basis for later econophysics-inspired order book models.

\subsection{Order-driven Market Modelling in Econophysics}

\subsubsection{The Order Book as a Reaction--Diffusion Model}
Bak, Paczuski and Shubik \cite{bak1997} proposed a reaction--diffusion model in which buy and
sell orders are treated as particles diffusing on a discrete price grid. A transaction occurs
when a buy and sell particle meet. While highly stylized, the model initiates the analogy
between order books and interacting particle systems. The resulting price dynamics exhibit
a Hurst exponent of $H=1/4$, far from empirical values.

\subsubsection{Introducing Market Orders}
Maslov \cite{maslov2000} extends the particle model by introducing fixed limit orders, market
orders that remove the best quote, and cancellation mechanisms. Orders are placed around
the prevailing price, and market orders generate price changes. This produces heavy-tailed
returns, though price dynamics remain subdiffusive ($H=1/4$).

\subsubsection{Deposition--Evaporation Models}
Challet and Stinchcombe \cite{challet2001} interpret the order book as an infinite one-dimensional
lattice with deposition (limit orders), annihilation (market orders), and evaporation 
(cancellation). The volatility behaviour is closer to empirical findings than earlier models
but still lacks over-diffusive properties.

\subsection{Empirical Zero-intelligence Models}

Mike and Farmer \cite{mikefarmer2008} represent a major advance in empirically calibrated
zero-intelligence modelling. Orders are placed according to Student-\(t\) distributions centered
around the best quotes. Cancellations depend on three empirically observed factors:
\begin{enumerate}
    \item Position in the order book: \(y(t) = \Delta(t)/\Delta(0)\);
    \item Imbalance between bid and ask queues: \(N_{\mathrm{imb}}(t)\);
    \item Total number of standing orders: \(N(t)\).
\end{enumerate}

The cancellation probability is given by:
\begin{equation}
P(C \mid y(t), N_{\mathrm{imb}}(t), N(t))
= 
A (1 - e^{-y(t)})
\bigl(N_{\mathrm{imb}}(t) + B \bigr)
\frac{1}{N(t)},
\label{eq:mike_farmer_cancellation}
\end{equation}
where \(A\) and \(B\) are fitted to data.

The model reproduces several important stylized facts (heavy tails, realistic spread
distribution), though it still encounters issues related to order-book stability and volatility
clustering.

% \subsection{Analytical Developments in Zero-intelligence Order Book Models}

% Smith et al.\ \cite{smith2003} study scaling properties of liquidity and price dynamics in a limit
% order book with Poissonian order flow. Mean-field approximations yield several analytical
% relations summarised in Table~\ref{tab:smith_scaling}.

% \begin{table}[h!]
% \centering
% \caption{Scaling relations from \cite{smith2003}. Here
% \(\epsilon := q/(\lambda_M / 2\lambda_C)\) is a granularity parameter,
% \(p_c := \lambda_M / (2\lambda_L)\) is a characteristic price interval, and
% \(f\) and \(g\) are slowly varying functions.}
% \label{tab:smith_scaling}
% \begin{tabular}{ll}
% \toprule
% \textbf{Quantity} & \textbf{Scaling relation} \\
% \midrule
% Average asymptotic depth & \( \lambda_L/\lambda_C \) \\
% Average spread & \( (\lambda_M/\lambda_L)\, f(\epsilon, \Delta P / p_c) \) \\
% Slope of depth profile & \( (\lambda_L^2 / \lambda_M \lambda_C)\, g(\epsilon, \Delta P/p_c) \) \\
% Short-term diffusion & \( \lambda_M^2 \lambda_C^{-1} \lambda_L^{-1} \epsilon^{-1/2} \) \\
% Long-term diffusion & \( \lambda_M^2 \lambda_C^{-1} \lambda_L^{-1} \epsilon^{1/2} \) \\
% \bottomrule
% \end{tabular}
% \end{table}

% Cont et al.\ \cite{cont2010} propose a simple but analytically tractable Markovian model in
% which limit orders arrive at distance \(i\) from the opposite best quote according to Poisson
% rates \(\lambda_i\), market orders arrive with rate \(\mu\), and standing orders cancel at intensity
% proportional to their queue size. All orders have unit size. The model admits a stationary
% distribution and yields closed-form expressions for various quantities of interest, including
% execution probabilities and mid-price movement probabilities.
\section{The Mathematical Structure of Zero-Intelligence Models}

This section develops a general mathematical framework for modelling limit order books
through Markovian, order-driven dynamics. Our objective is to understand how the
structural features of the book, namely, its depth profile, order-flow composition, and quote
adjustments, which govern the evolution of market observables such as price and spread.  

We begin with a minimalistic agent-based representation of a Bachelier-type price process,
which illustrates how classical diffusion models arise from microstructural order flows.
We then introduce the full Markovian description of a limit order book driven by
independent Poisson processes, including the evolution of queue sizes, best quotes, and
price increments.

% ----------------------------------------------------------------------
\subsection{An Elementary Approximation: Perfect Market Making}
% ----------------------------------------------------------------------

To motivate the Markovian modelling of limit order books, consider a highly stylised
``perfect market making'' framework. Although simple, it reveals how macroscopic
diffusion-type price behaviour emerges from discrete order-flow events.

We impose the following assumptions:
\begin{itemize}
    \item The order book initially contains one unit-sized limit order at every price level above
    the best ask \(P_A(0)\) and below the best bid \(P_B(0)\). The initial spread is one tick.
    \item Buy and sell market orders arrive as independent Poisson processes
    \( M^+(t) \) and \( M^-(t) \)
    with respective intensities \( \lambda_+ \) and \( \lambda_- \).
    \item A single liquidity provider maintains a one-tick spread.  
    After a buy (respectively sell) market order, they submit a buy (respectively sell) limit order
    with probability \(u\), and an order on the opposite side with probability \(1-u\).
\end{itemize}

The resulting mid-price process satisfies:
\begin{equation}
dP(t)
= \Delta P \bigl( dM^+(t) - dM^-(t) \bigr) Z,
\end{equation}
where \(Z\) is a Bernoulli random variable with: 
\[
\mathbb{P}(Z=1)=u, \qquad \mathbb{P}(Z=0)=1-u,
\]
and \( \Delta P \) denotes one tick.

The infinitesimal generator is:
\begin{equation}
Lf(P)
=
u\lambda_+ \bigl( f(P+\Delta P) - f(P) \bigr)
+
u\lambda_- \bigl( f(P-\Delta P) - f(P) \bigr).
\end{equation}

Expanding in Taylor series and grouping first-order and second-order terms yields  
\begin{align}
Lf(P)
&=
\frac{u(\lambda_+ + \lambda_-)(\Delta P)^2}{2}\,
\frac{\partial^2 f}{\partial P^2}
+
u(\lambda_+ - \lambda_-)\Delta P\,
\frac{\partial f}{\partial P}
+ o((\Delta P)^2).
\end{align}

Under the scaling limits:
\begin{equation}
u(\lambda_+ + \lambda_-)(\Delta P)^2 \to \sigma^2,
\qquad
u(\lambda_+ - \lambda_-)\Delta P \to \mu,
\end{equation}
the generator converges to the diffusion operator:  
\begin{equation}
Lf(P)
=
\frac{\sigma^2}{2}\frac{\partial^2 f}{\partial P^2}
+
\mu \frac{\partial f}{\partial P},
\end{equation}
corresponding to a Wiener process with drift.  
This example illustrates how classical continuous-time models may arise as scaling limits
of simple order-driven mechanisms.

% ----------------------------------------------------------------------
\subsection{Limit Order Book Representation and Notation}
% ----------------------------------------------------------------------

We now introduce the general stochastic framework for a limit order book driven by
Poissonian order flows. The book is represented in a moving frame of fixed size \(K\) on
each side of the spread.

Let
\[
(a(t);\, b(t)) =
(a_1(t),\dots,a_K(t);\; b_1(t),\dots,b_K(t)),
\]
where
\begin{itemize}
    \item \(a_i(t)\) is the queue size \(i\) ticks above the best bid,
    \item \(b_i(t)\) is the queue size \(i\) ticks below the best ask.
\end{itemize}

Queue sizes live in \(q\mathbb{Z}\), where \(q\) is the lot size.  
Define cumulative depths:
\begin{equation}
A_i = \sum_{k=1}^i a_k,
\qquad
B_i = \sum_{k=1}^i |b_k|.
\end{equation}

The indices of the first non-empty ask and bid queues are given by the generalized inverse:
\begin{equation}
A^{-1}(q') := \inf\{p : A_p > q'\},
\qquad
B^{-1}(q') := \inf\{p : B_p > q'\}.
\end{equation}

The spread in ticks is therefore
\begin{equation}
i_S = A^{-1}(0) = B^{-1}(0).
\end{equation}

% --------------------------------------------
\paragraph{Boundary Conditions.}
To ensure that best quotes remain well defined and the state process remains Markovian,
we impose constant boundary depths \(a_\infty\) and \(b_\infty\) at price levels leaving the
moving frame. This prevents the book from emptying and keeps the spread bounded.

% ----------------------------------------------------------------------
\subsection{Order Arrivals and State Transitions}
% ----------------------------------------------------------------------

The evolution of the order book is driven by three types of events:
\begin{itemize}
    \item Market orders, arriving as Poisson processes \(M^\pm(t)\) with intensities
    \( \lambda_M^\pm \);
    \item Limit orders at level \(i\), arriving with intensities \( \lambda_{L,i}^\pm \);
    \item Cancellations at level \(i\), occurring with intensities
    \( \lambda_{C,i}^\pm a_i \) and \( \lambda_{C,i}^\pm |b_i| \).
\end{itemize}

All orders have fixed size \(q\).  
Limit orders arrive relative to the best opposite quote on their respective sides.

The queue sizes satisfy stochastic difference equations combining:
\begin{itemize}
    \item Increases from limit-order arrivals,
    \item Decreases from cancellations and market-order executions,
    \item \emph{Shift operators} that reindex queues when a best quote is consumed or
    a limit order arrives inside the spread.
\end{itemize}

These shift operators maintain exactly \(K\) visible levels in the moving frame and ensure a
consistent Markovian representation.

% ----------------------------------------------------------------------
\subsection{Price Dynamics}
% ----------------------------------------------------------------------

Let \(P_A(t)\) and \(P_B(t)\) denote the best ask and bid prices.  
Events at or within the spread affect these quotes:

\begin{itemize}
    \item A market buy consumes depth on the ask side,
    \item A limit order placed inside the spread improves a quote,
    \item A cancellation at the best queue may widen the spread.
\end{itemize}

The best ask satisfies:
\begin{align}
dP_A(t)
= \Delta P \Bigl[
&\bigl(A^{-1}(q)-i_S\bigr)\, dM^+(t)
-
\sum_{i=1}^K (i_S - i)^{+}\, dL^+_i(t)
\\
&\qquad
+\bigl(A^{-1}(q)-i_S\bigr) dC^+_{i_S}(t)
\Bigr].
\end{align}

A symmetric expression governs \(P_B(t)\).  
The mid-price \(P(t) = (P_A(t)+P_B(t))/2\) and spread \(S(t)=P_A(t)-P_B(t)\) follow:

\begin{align}
dP(t)
&= 
\frac{\Delta P}{2}
\Bigl[
\bigl(A^{-1}(q)-i_S\bigr)\, dM^+(t)
-
\bigl(B^{-1}(q)-i_S\bigr)\, dM^-(t)
\Bigr]
\notag\\[4pt]
&\quad
-
\frac{\Delta P}{2}
\sum_{i=1}^K (i_S - i)^{+}\, dL^+_i(t)
+
\frac{\Delta P}{2}
\sum_{i=1}^K (i_S - i)^{+}\, dL^-_i(t)
\notag\\[4pt]
&\quad
+
\frac{\Delta P}{2}
\bigl(A^{-1}(q)-i_S \bigr)\, dC^+_{i_S}(t)
-
\frac{\Delta P}{2}
\bigl(B^{-1}(q)-i_S \bigr)\, dC^-_{i_S}(t).
\end{align}

Similarly, the spread dynamics satisfy
\begin{align}
dS(t)
&=
\Delta P
\Bigl[
\bigl(A^{-1}(q)-i_S\bigr)\, dM^+(t)
+
\bigl(B^{-1}(q)-i_S\bigr)\, dM^-(t)
\Bigr]
\notag\\[4pt]
&\quad
-
\Delta P
\sum_{i=1}^K (i_S - i)^{+}\, dL^+_i(t)
-
\Delta P
\sum_{i=1}^K (i_S - i)^{+}\, dL^-_i(t)
\notag\\[4pt]
&\quad
+
\Delta P
\bigl(A^{-1}(q)-i_S \bigr)\, dC^+_{i_S}(t)
+
\Delta P
\bigl(B^{-1}(q)-i_S \bigr)\, dC^-_{i_S}(t).
\end{align}

These expressions provide explicit links between the local liquidity profile of the book and
the resulting dynamics of price and spread. Depth at the best levels determines the
immediate impact of market orders, while the limit-order flow within the spread governs
quote adjustments and spread behaviour.

\section{Conclusion}

These insights form the foundation for the next chapter, where we move from theory to
practice by simulating limit order book dynamics under the modelling assumptions
introduced here. Through numerical experimentation, we will investigate how the
interaction of arrivals, cancellations, and liquidity profiles gives rise to the statistical
properties observed in real markets.

